\slajdid{02-s045}
\begin{frame}[label=02-s045]
\gusttekst
Za realizaciju u obliku zbira proizvoda
\[
F=\bar D\bar B+D\bar C\bar A
\]
potrebno je 1 dvoulazno I kolo, 1 troulazno I kolo i 1 dvoulazno ILI kola, odnosno ukupan broj potrebnih ulaza je 7. Za realizaciju u obliku proizvoda zbirova
\[
F=(D+\bar B)(\bar D+\bar C)(\bar D+\bar A)
\]
potrebno je 3 dvoulazna ILI kola i 1 troulazno I kolo, odnosno ukupan broj potrebnih ulaza je 9. Sa stanovišta klasične minimizacije je minimalnija realizacija u obliku zbira proizvoda.

\alert{U startu se ne zna koja je realizacija minimalnija. Treba izvesti obe pa na osnovu njihovih minimalnih izraza izvesti zaključak šta je zapravo minimalna realizacija.}
\end{frame}
